% ---
% title: Handwritten Medicine MCQ Answers
% date: 2026-09
% category:
% nodes:
% links:
% unlisted: true
% ---
\documentclass{article}
\usepackage[utf8]{inputenc}
\usepackage{amsmath, amssymb, amsthm}
\usepackage{geometry}
\usepackage{graphicx}
\usepackage{hyperref}

\geometry{a4paper, margin=1in}

\title{Handwritten Medicine MCQ Answers}
\date{September 2026}
\author{Daksh Mehta}

\begin{document}

Answers to the questions on the \href{handwritten-medicine-mcqs.html}{main MCQ page}. 

\tableofcontents

\section{Hypersensitivity reactions in SLE}

\textbf{Answer: b. Type II}

The vignette describes antiphospholipid syndrome, which is secondary to SLE here. The antibodies against beta 2-glycoprotein I and cardiolipin bind their target antigens where those antigens sit on cell surfaces, most importantly on platelets and vascular endothelium. Binding activates those cells directly, producing the the prothrombotic state responsible for her deep vein thrombosis and her recurrent miscarriages. Antibody directed against a cell-surface antigen is a type II reaction.

The prolonged activated partial thromboplastin time is a laboratory artefact rather than evidence of bleeding risk. The antibodies interfere with the phospholipid reagent used in the assay, so the number goes up while the patient clots more, not less. A normal bleeding time alongside it is the clue that this is not a genuine coagulopathy.

Type III is the tempting distractor, because SLE itself is the textbook immune complex disease and is responsible for the glomerulonephritis, arthritis and malar rash seen elsewhere in the condition. The question asks specifically about the mechanism behind the current presentation, and thrombosis in antiphospholipid syndrome is not caused by complex deposition. Type I would require IgE and mast cell degranulation (think typical allergic reaction), and type IV would be T cell mediated and delayed (think contact dermatitis) neither of which fits.

\end{document}
